\section{Method}
\label{sec:method}

\subsection{Scene and renderer}
The canonical scene is a procedural grass field: a jittered grid of blades on a
ground plane, with a few coloured spheres, lit by a single directional sun over
a gradient sky. Each blade is a canonical shape in local space (four triangles
at full detail) placed by an affine transform $x \mapsto Ax + b$ encoding
position, yaw, tilt, and height. Everything is deterministic given a seed. Our
rendered scene uses \SceneBlades\ blades (\SceneTris\ triangles) at
$64\times48$; the accompanying \texttt{gpu\_full} config scales this up without
code changes.

The renderer is a Monte-Carlo path tracer: a primary ray per pixel (jittered for
anti-aliasing), a shadow ray to the sun, and one cosine-weighted diffuse bounce
whose contribution is estimated from direct light at the second hit or the sky.
Averaging $N$ samples per pixel (spp) reduces variance. A deterministic centre
ray also records a G-buffer (albedo, normal, depth) used by the denoiser.

\subsection{Cost as a currency}
We report cost as the count of \emph{ray/primitive intersection tests}
(triangle, sphere, plane, and bounding-box tests), incremented by every routine
that performs one. Unlike wall-clock time, this is exact and machine
independent. For the acceleration study we count the tests to cast one primary
plus one shadow ray per pixel; for the Pareto study we use
\begin{equation}
\text{cost}(\text{method}, N) \;=\; \text{tests}_{\text{1 pass}} \times N,
\end{equation}
a single currency that combines geometric cost (the intersection structure) with
sampling cost (spp), normalised so brute force at the reference sample count is
$1$.

\subsection{The four techniques}
\paragraph{Brute force (baseline).} Every ray is tested against every triangle:
$O(\text{rays}\times\text{triangles})$.

\paragraph{BVH.} A median-split bounding-volume hierarchy over triangle
centroids. A ray descends the tree, testing a box before its subtree and only
reaching the few triangles in leaves it cannot rule out. The image is identical
to brute force; only the test count changes.

\paragraph{Instancing.} A two-level BVH: the canonical blade is stored once as a
BLAS, and each blade is an instance in a TLAS. A ray is transformed into an
instance's local frame ($A^{-1}(o-b)$, $A^{-1}d$, without renormalising $d$ so
the hit distance $t$ is preserved) and tested against the shared blade. Geometry
is stored once rather than duplicated $K$ times; the image is bit-identical.

\paragraph{Level of detail.} Beyond a distance threshold, a blade is baked with a
cheaper shape --- a single triangle (\emph{simplified}) or a flat two-triangle
card (\emph{billboard}). The threshold in $[0,1]$ is the fraction of the
camera-distance range kept at full detail; lower means cheaper and more
approximate.

\paragraph{Sample budget and denoising.} We sweep spp and, at each, apply a
cross-bilateral denoiser whose weights combine spatial distance, colour
similarity, and geometric similarity (normal agreement to the eighth power and a
depth term), so it smooths noise within a surface without blurring across
silhouettes.

\subsection{Metrics}
Quality is measured against a high-sample reference (\RefSpp\ spp) with PSNR and
windowed SSIM \cite{wang2004image}, both implemented in numpy. Lossless methods
attain SSIM $=1$ by construction (verified to floating-point in the tests);
lossy methods produce a curve.
